% Section 1.3, from lectures/L01/README.md: the objectives, the evaluation questions, and the next
% lecture.

\section{Review}
\label{c1:sec:review}

\needspace{6\baselineskip}
\subsection*{What you should be able to do now}
After this chapter, you should be able to:
\begin{itemize}
  \item \textbf{Read the provided register-map package} (\code{uart_def}) and explain why its
    \code{natural} bit-\textbf{position} constants must match the protocol, and why the C++
    \file{register_map.hpp} agrees with it value for value wherever a name exists on both sides.
  \item \textbf{Build the \code{uart_top} skeleton}, meaning the entity, the provided
    \code{reset_sync}, the provided transport, and the signals those three need, so that it analyzes
    cleanly with no datapath blocks present.
  \item \textbf{Derive the internal signals from the entities they connect}, rather than from a
    list, and apply the \code{spi_} prefix and \code{_s2} suffix conventions the provided modules
    already use.
  \item \textbf{Explain top-down structural design}: why fixing the top first fixes each block's
    port contract, and why the system testbench is gated, skipped rather than failed, until
    \chapref{5}.
\end{itemize}

\needspace{6\baselineskip}
\subsection*{Questions to test yourself}
You should be able to explain:
\begin{enumerate}
  \item Why build \code{uart_top} before any block it instantiates, and what does fixing the top
    first fix about every datapath module you write later?
  \item The reset is asserted asynchronously but released synchronously: what can go wrong if the
    release is not synchronized, and why can \code{sync} (\chapref{3}) not do this job?
  \item \code{uart_top_tb} binds \code{uart_top}'s ports positionally: what class of wiring mistake
    does that leave the analyzer unable to catch, and at which point in the book does it finally
    surface?
\end{enumerate}

\needspace{6\baselineskip}
\subsection*{Looking ahead}
The datapath begins in \chapref{2}: \code{baud_gen}, the shared time base, and \code{uart_tx}, the
transmitter, each verified against its own testbench and instantiated into the \code{uart_top} you
built here.
